% !TeX root = ../教材题库_全部条目_转写工作稿.tex
% Source: XJ-B1-C05-RV-P206-Y1435
% 第5章 / 章末复习 复习题五 / 教材 p.206 / PDF 第213页
% OCR未校对：公式与文字均可能有误；未完成数学TeX转写；禁止直接发给学生
% 原题及相邻原解截图（核对用）：../题目原图/XJ-B1-C05-RV-P206-Y1435-01.webp
% 本文件只在转写完成后作为题干使用；OCR 条目尚不是可靠题干。
\begin{bankitem}{复习题五 第8题}{第5章 / 章末复习 复习题五 / 教材 p.206 / PDF 第213页}{XJ-B1-C05-RV-P206-Y1435}{OCR未校对：公式与文字均可能有误；未完成数学TeX转写；禁止直接发给学生}
\begin{ocrdraft}
[教材p.206 / PDF第213页]
\par
8．利用“五点法”作出函数一sin产e@，2和函数3，=co丷r，xe@，
\par
2冠的图象，并指出一sinx是减函数且，x是增函数时的取值范围．
\end{ocrdraft}
\end{bankitem}
